\documentclass[11pt]{article}

\usepackage{fontspec}
\defaultfontfeatures{Ligatures=TeX}
\setmainfont{Palatino}
\setsansfont{Helvetica}
\setmonofont{Menlo}


\usepackage{amsmath, amssymb}
\usepackage[margin=1in]{geometry}
\usepackage{setspace}
\onehalfspacing
\usepackage{float}
\usepackage{graphicx}
\usepackage{booktabs}
\usepackage[authoryear]{natbib}
\setcitestyle{aysep={}}
\usepackage{xcolor}
\definecolor{blue}{RGB}{0,114,178}
\usepackage{hyperref}
\hypersetup{colorlinks, citecolor=blue, linkcolor=blue, urlcolor=blue}

\begin{document}

\title{When Conventions Collapse: Standing Committee Chair Allocation in the Korean National Assembly under Supermajority Conditions}
\author{KNA Research Agents (AI-generated)\thanks{Experimental output - kna-research-agents.com. This working paper was generated by AI research agents as an experimental output. It has not been peer-reviewed or fact-checked. Do not cite or use in any academic, policy, or professional context. For external double-blind submission, the title page would be anonymized.}}
\date{\today}
\maketitle
\thispagestyle{empty}

\begin{abstract}
\noindent I examine the durability of informal conventions that bind a legislative majority to share procedural power with the minority. The focus is on four conventions governing standing committee chair allocation in the Korean National Assembly, in particular the long-standing rule that the Legislation and Judiciary Committee chair belongs to a party other than the Speaker's. Using a hand-coded panel of 149 committee-half cells across the 18th through 22nd Assemblies (2008-2026), I document that three of the four conventions identified by \citet{jung2018} held cleanly through the 18th-20th Assemblies and broke during the 22nd, when a single bloc held a near-supermajority. The aggregate cooperative legislative-productivity pattern that \citet{yunan2018} attributed to the post-2012 reform era is consistent with a sign reversal in the same period, although the design admits a composition explanation that the within-committee fixed-effects analysis only partially rules out. The findings suggest that residual chair-allocation norms may function as strategic precommitment devices that lose their stabilizing role when alternation expectations recede.
\end{abstract}

\bigskip
\noindent\textbf{Keywords:} committee chairs, legislative institutions, Korean National Assembly, supermajority, agenda control

\bigskip
\clearpage

\section{Introduction}

Procedural conventions in legislatures often outlive both the formal rules that produced them and the partisan circumstances that did. The convention by which the Korean National Assembly's Legislation and Judiciary Committee (법제사법위원회) is chaired by a member of a party other than the one to which the Speaker belongs is one such convention. It has no statutory basis in the National Assembly Act of Korea. It is not enumerated among the supermajority-protective rules introduced by the 2012 National Assembly Advancement Act, which \citet{jeong2023} treats as the formal codification of cooperation requirements between the largest and second-largest blocs. Yet for at least the three Assemblies preceding the 21st (2020-2024), the convention bound even sizable majorities to cede the gavel of the chamber's most consequential procedural committee to the opposition.

The convention broke. In the first half of the 21st Assembly the Democratic Party of Korea (더불어민주당; hereafter Democratic Party or DPK), having won a 174-seat majority in the April 2020 general election and holding the Speakership, took the Legislation and Judiciary Committee chair as well. The convention was partially restored in the second half of that Assembly through inter-party negotiation following the 2022 presidential transition, when the executive moved to the People Power Party (국민의힘; hereafter PPP) under President 윤석열. In the 22nd Assembly, after the Democratic-led coalition reached approximately 175 seats in the April 2024 general election, the convention collapsed completely. Three of the four chair-allocation rules that \citet{jung2018} identified as stable institutional conventions failed simultaneously in the first half of the 22nd Assembly. Only the chair of the National Defense Committee continued to follow its conventional party assignment. The PPP retained the National Defense chair against an opposition supermajority that controlled the rest of the chair-share negotiation.

The puzzle this paper addresses is twofold. First, why did a convention that survived three earlier moments of single-party legislative dominance, including the 19th Assembly's conservative majority and the 20th Assembly's progressive transition, fail when the Democratic Party crossed the threshold into supermajority territory? Second, what happens to legislative productivity when the cross-party allocation convention that \citet{yunan2018} credited with sustaining the post-2012 cooperative pattern is no longer in force? The strategic-precommitment framework developed by \citet{jeong2023} for the 2012 reform suggests that residual chair-allocation norms function as informal insurance against future minority status; when alternation expectations recede, the norm's strategic logic collapses, and the norm is rationally abandoned. The Korean political-science literature on chair selection, which I review in the next section, has carefully mapped which legislators become chairs \citep{kang2023, kang2024}, but it has not asked what changes when the entire allocation regime breaks down.

I investigate these questions using a panel of 149 committee-half cells covering all 17 standing committees across the 18th through 22nd Assemblies, where the cell count reflects the 153 potential committee-half slots minus four cells dropped because the relevant committee was either restructured or had no member-introduced bill activity in the half. The chair-tenure information is reconstructed from the kr-hearings-data corpus, in which the role of presiding chair is tagged at the speech level for every committee meeting. Modal-presider classification within each committee-half yields chair attributions for 150 chair spells, with mid-term rotation observed in 62 of 64 within-Assembly transitions where both halves are observed and the committee structure is held fixed (96.9 percent), consistent with the two-year norm enshrined in National Assembly Act §41. Bill-level legislative productivity is taken from the master\_bills corpus across the same period.

The paper makes three empirical contributions. First, I document that the cross-party 법사위 convention, the ruling-party 운영위 convention, the ruling-party 행안위 convention, and the ruling-party 국방위 convention identified by \citet{jung2018} held in 31 of 32 applicable cells in the 18th through 20th Assemblies and broke in 4 of 8 applicable cells in the 21st through 22nd Assemblies, with the breaks concentrated in the supermajority-defined 22nd Assembly's first half. Second, I show that the legislative-productivity pattern is consistent with a sign reversal predicted by the strategic-precommitment framework: pre-supermajority opposition-chaired committees outperformed ruling-chaired committees in passage rate by a substantively large margin, while post-supermajority ruling-chaired committees now outperform opposition-chaired committees. Third, I report a within-Assembly placebo from the 21st Assembly's mid-term restoration of the judiciary convention, in which bill volume contracted substantially and the passage rate fell when the convention was restored.

The paper proceeds in six parts. Section 2 reviews the Korean institutional literature on chair allocation, the comparative literature on parliamentary committee chairs, and the strategic-precommitment framework that motivates the supermajority hypothesis. Section 3 describes the data and the modal-presider classification. Section 4 presents the convention-status panel, the productivity sign-reversal pattern, and a within-committee fixed-effects specification that probes whether the reversal is mechanism or composition. Section 5 discusses the bundled-negotiation alternative explanation that \citet{leekim2022} document, the polarization alternative that I cannot rule out at the available sample size, and the implications for the d'Hondt mechanical-allocation reform proposal. Section 6 concludes with the limitations of an N=2 supermajority panel and the back-extension to the 13th-17th Assemblies that the project's next stage will pursue.

The contribution is descriptive in the strict sense that no causal estimate of the supermajority effect is possible at the available number of treated Assemblies. The paper's value lies in establishing the empirical pattern, in placing it within the strategic-precommitment apparatus that \citet{jeong2023} developed for the Korean case, and in distinguishing what the data can support from what they cannot. I treat this descriptive-only stance as the primary methodological commitment of the paper.

\section{Literature and Theory}

The Korean political-science literature on standing committee chair allocation has produced a mature selection-side body of work. \citet{kang2023}, in the most direct precedent for the present paper, examines which legislators become standing committee chairs and committee secretaries in the 20th National Assembly. He finds that party loyalty, measured by within-party voting cohesion, predicts leadership selection in the minority party but not in the majority party, while ideological distance from the party median predicts non-selection again only in the minority bloc. His interpretation is that the minority party uses chair and secretary posts as the primary reward currency for party loyalty in a way that the majority party does not. \citet{kang2024} extends the framework to committee assignment more broadly and documents that party loyalty plus prior assignment experience predict assignment to committees that confer reelection advantage, while individual-level electoral safety is not a significant predictor.

\citet{choikoo2018} provide the theoretical anchor for the Korean case. They critically test the three US-derived theories of committee composition (distributive, informational, partisan) on the 18th and 19th Assemblies and conclude that the partisan theory of \citet[ch.~1-3]{coxmccubbins2005} dominates the alternatives. Committee composition reflects the party's ideological direction, and the alignment is stronger in the second half of each Assembly term than the first. Their finding is consistent with \citet{kang2024}'s assignment-experience-persistence result and with the party-cartel framework that motivates the present paper's interpretation of chair allocation as a partisan-control technology rather than an informational-specialization technology in the sense of \citet{gilligankrehbiel1987}.

The institutional anchor is \citet{leekim2022}. Surveying the 13th through 21st Assembly chamber-formation negotiations, they document three recurring pathologies: post-election member-poaching attempts, chair-share negotiations bundled with unrelated bargaining agendas (such as prosecutorial reform or media law), and repeated conflicts over chair-share ratios that delay 원구성 (chamber formation) by 30 to 90 days each Assembly. Their reform proposal is the application of the d'Hondt highest-average method to chair allocation, which would automate the share-ratio and choice-ordering decisions that bundled negotiation otherwise distorts. For the present paper, \citet{leekim2022} supplies both the empirical typology of negotiation pathologies and the policy implication: if the convention regime has collapsed, mechanical allocation rules become the only credible commitment device.

\citet{jung2018} is the descriptive anchor. His survey of 13th through 21st Assembly leadership formation documents two stable conventions that are not enshrined in statute. First, the Speaker's office increasingly goes to the largest party, with non-largest-party choice acting as a swing variable when no clear majority exists. Second, the Legislation and Judiciary Committee chair is allocated to a party other than the Speaker's regardless of the seat ratio. The second convention is the central object of the present paper. \citet{jung2018} treats it as a stable convention without offering a strategic-precommitment account; the present paper supplies that account by drawing on \citet{jeong2023}.

\subsection{Strategic Precommitment as the Theoretical Motor}

\citet{jeong2023} argues that the 2012 National Assembly Advancement Act was adopted by the majority party to shift the legislature from a norm-based to a rule-based regime because the prior norm regime produced chaos and uncertainty in lawmaking. The strategic motor is alternation: a majority that anticipates becoming a minority in the near future has reason to support institutional protections that bind future majorities, because today's protection is tomorrow's insurance. The 2012 reform was a formal-rule substitute for failing informal norms in the 17th-19th Assembly period, when alternation expectations were high.

The framework yields a prediction structure for residual norms not codified by the 2012 reform. The cross-party 법사위 convention is one such residual norm. It is functional under conditions where today's majority anticipates being tomorrow's minority. When the alternation-fear mechanism is disabled, either by formal rule or by structural circumstance such as a stable supermajority that does not anticipate near-term alternation, the strategic logic that sustained the residual norm collapses, and the norm is rationally abandoned.

The 2020 and 2024 general elections produced Democratic-led coalitions of approximately 174 seats (2020) and 175 seats (2024) in a 300-seat chamber. The "approximately" qualifier reflects the inclusion of small allied parties and one or two independents whose caucus alignment fluctuated; using a Democratic-Party-only count yields 174 and 171 respectively. The post-2012 supermajority threshold (180 seats for cloture-equivalent action under the National Assembly Advancement Act) was approached but not crossed in either election. The convention break I document in the empirical sections occurred at seat shares below the formal supermajority threshold, suggesting that the strategic-precommitment mechanism may operate at a lower seat-share threshold than the formal cloture-equivalent rule requires.

\subsection{The Comparative Parliamentary Anchor}

\citet{fmv2017} is the contemporary parliamentary-systems anchor. They argue that committee chairs in coalition parliamentary systems function as policy-review delegates, with chairs from coalition-partner parties using the post to monitor cabinet ministers from rival parties. The theoretical move is to treat the chair as a bureaucratic-control technology rather than a pure agenda-setting technology. The framework yields a prediction for the Korean case: if the cross-party 법사위 chair functioned as a \citet{fmv2017} review delegate when the convention held, then its committee should have shown elevated bill-amendment rates and lower direct-passage rates relative to majority-party-chaired committees, conditional on bill content. The convention-break period offers a test of whether the review-delegate function survives the loss of the convention.

The Cox-McCubbins party-cartel framework \citep[ch.~3]{coxmccubbins2005} interprets committee assignments and chair allocations as instruments of negative agenda control deployed by the majority leadership to protect the party brand. The framework predicts that majority parties retain control of high-leverage committees and concede low-leverage committees as a function of internal coalition management. The Korean cross-party 법사위 convention runs against this prediction: the majority party voluntarily ceded the highest-leverage procedural committee to the opposition for at least three Assemblies. The strategic-precommitment framework I draw from \citet{jeong2023} provides one resolution: the cross-party concession was an insurance investment, not a coalition-management concession. \citet{gilligankrehbiel1987} would offer a different resolution under their informational rationale, but \citet{choikoo2018}'s adjudication of the three US theories on Korean data found that the informational interpretation does not fit, leaving the strategic-precommitment reading as the dominant available account.

\citet{mickler2017} provides the cross-national counter-bench. He analyzes 1990-2013 Bundestag committee assignments and finds that constituency demands, regional factions, reassignment patterns, and occupational background all matter for assignment. Crucially, the German Bundestag does not allocate committee assignments through a single proportional rule. The cross-party 법사위원장 convention that the Korean Assembly developed was a stricter constraint than mature European parliamentary practice, which means the pre-2020 Korean convention was a high-effort coordinative equilibrium rather than the default cross-national pattern. The Bundestag bench is therefore a counter-example rather than a comparative target. I cite \citet{mickler2017} in this paper for the institutional-comparison purpose only.

\subsection{Selection versus Behavior, and the Carey-Shugart Pipeline}

The selection literature surveyed above is cross-sectional. It asks which legislators become chairs, given the chair-allocation regime in place. The present paper asks a different question. It asks what happens to the chair-allocation regime itself, and what consequences regime change has for legislative productivity. The two questions are linked through \citet{kang2023}'s minority-party-loyalty finding: the selection mechanism conditional on regime is not the same in the majority and minority blocs, which means that a regime change altering which bloc holds which chairs should propagate into changes in the type of legislator selected, conditional on bloc.

\citet{careyshugart1995} provides the broader framework for understanding why selection mechanisms differ across electoral and party-organizational contexts. Korean legislators face a mixed-member majoritarian system in which seniority and party loyalty are weighted differently in proportional-representation and single-member-district pathways. \citet{kang2024}'s assignment-experience-persistence finding suggests that the selection rule rewards repeat assignment, which biases the selection pool toward members who have already accumulated committee-specific human capital. The behavioral consequence I trace in the empirical sections, in which the convention-break period is associated with a productivity sign reversal, is consistent with a selection-based interpretation: when the cross-party convention breaks and the majority party appropriates committees it would not otherwise hold, the selection pool for those committees shifts in ways that change behavior.

\subsection{Hypotheses}

Three testable expectations follow from the literature. First, the convention-durability expectation. Below the supermajority threshold and with active alternation expectations, the four \citet{jung2018} conventions should hold in close to all applicable cells; the 18th through 20th Assemblies satisfy these conditions and provide the unbroken-convention benchmark (H1). Second, an expectation about supermajority-induced break, consistent with the single observed near-supermajority case in the panel rather than a generalization across many: under conditions of opposition supermajority approaching but not exceeding the formal cloture-equivalent threshold of about 175 seats out of 300, residual norms that lack statutory protection should collapse selectively, and the 22nd Assembly's first half is the case in which I observe such collapse (H2). Third, the productivity-reversal expectation. If the cooperative-legislation pattern that \citet{yunan2018} documented for the 19th Assembly is a behavioral consequence of the convention regime rather than an autonomous effect of the 2012 formal rule, then breakdown of the regime should be accompanied by sign reversal or substantial attenuation in the chair-party-passage relationship (H3).

These three expectations are not independent. H1 and H2 jointly constitute the convention-break finding; H3 is the behavioral consequence of the regime change H2 documents. I treat H1 and H2 as the descriptive headline and H3 as the mechanism inquiry. Per the project's small-N guardrail, H2 is best read as an expectation consistent with the single near-supermajority case observed (the 22nd Assembly's first half, with the 21st Assembly's first half as a partial precursor), not as an inferential generalization across many treated units.

\section{Data and Method}

\subsection{Data}

The chair-tenure dataset covers all 17 standing committees of the Korean National Assembly across the 18th through 22nd Assemblies (May 2008 through April 2026). Chair attribution is derived from the kr-hearings-data corpus, in which the role of presiding chair is tagged at the speech level for every committee meeting. I restrict to records with hearing\_type set to standing committee (상임위원회) and role set to chair, then assign each (term, committee, half-of-term) cell to the legislator who delivered the largest number of chair-tagged speeches in that cell. The half-of-term split corresponds to the two-year rotation rule established by National Assembly Act §41, under which standing committee chairs serve two-year terms within a four-year Assembly. The 22nd Assembly's second half had not begun at the time of writing and is therefore absent from the panel.

The bill-level data are drawn from the master\_bills corpus across the 18th through 22nd Assemblies. Each bill carries a sponsoring legislator, a sponsoring party, a referring committee, a proposal date, and a binary passed indicator. I restrict the bill panel to member-introduced bills (의원 발의) and exclude government bills and committee-alternative bills. Committee assignment of bills follows the master\_bills committee\_nm field; half-of-term assignment of bills follows the proposal-date field relative to the term-start date.

Table~\ref{tab:desc} reports descriptive statistics for the chair-tenure panel and the bill-level panel. The panel contains 17 standing committees observed across nine half-Assemblies (18-H1 through 22-H1), yielding 153 potential committee-half cells. Three cells are dropped because the relevant committee was restructured during the half so that no single committee-name unit existed for the full half, leaving 150 chair-attributable spells. Of these, 149 cells have at least one member-introduced bill and constitute the bill-level analysis panel. The 36 \citet{jung2018}-applicable cells correspond to the four convention-bearing committees (Legislation and Judiciary, House Steering, Public Administration and Security, National Defense) observed across nine halves. The mid-term rotation rate is computed over the 64 within-Assembly chair-transitions in which both halves of an Assembly are observed and the committee identity is held constant across halves.

\begin{table}[ht]
\centering
\caption{Descriptive Statistics, 18th-22nd National Assembly Panel}
\label{tab:desc}
\begin{tabular}{lcccc}
\toprule
Variable & N & Mean & SD & Range \\
\midrule
Chair spells (committee-half cells) & 150 & n/a & n/a & 18--22 NA \\
Jung-applicable convention cells & 36 & n/a & n/a & 4 conventions \\
Mid-term chair-rotation rate & 64 & 0.969 & n/a & n/a \\
Member-introduced bills & 85,423 & n/a & n/a & 18--22 NA \\
Bill passage rate (overall) & 85,423 & 0.273 & 0.446 & 0--1 \\
Bills per chair-half cell & 149 & 573 & 412 & 18--2,047 \\
High-stakes committee bills (share) & 22,668 & 0.265 & n/a & n/a \\
\bottomrule
\multicolumn{5}{l}{\footnotesize Chair attribution by modal-presider rule on chair-tagged speeches.} \\
\multicolumn{5}{l}{\footnotesize n/a = not applicable for count statistics or single-cell denominators.} \\
\end{tabular}
\end{table}

The high-stakes committee classification follows the topic-gate pre-registration: high-stakes is the set of five committees comprising the Special Committee on Budget and Accounts, the Legislation and Judiciary Committee, the House Steering Committee, the National Policy Committee, and the Strategy and Finance Committee. Low-stakes is the residual 12 standing committees. The high-stakes set processes between 23 and 27 percent of bill volume across the panel, which is substantively meaningful without being so large as to drown out comparison.

\subsection{Identification Strategy}

The primary specification is a two-way fixed-effects regression of bill-level passage on chair-party indicators interacted with a post-supermajority indicator, with committee and half fixed effects:

\begin{equation}
\text{Passed}_{ict} = \beta_1 \text{ChairRuling}_{ct} + \beta_2 \text{PostSuper}_t + \beta_3 (\text{ChairRuling}_{ct} \times \text{PostSuper}_t) + \delta_c + \mu_h + \epsilon_{ict}
\label{eq:main}
\end{equation}

where $i$ indexes bills, $c$ indexes committees, and $t$ indexes Assembly-halves. ChairRuling is an indicator that the modal presider for the (committee, half) cell belongs to the same party as the executive (the President's party). PostSuper is an indicator for the 21st and 22nd Assemblies, the period in which the opposition Democratic-led coalition approached the supermajority threshold. The committee fixed effect $\delta_c$ absorbs time-invariant differences in passage propensity across the 17 standing committees; the half fixed effect $\mu_h$ absorbs Assembly-half-level shocks common to all committees. Standard errors are clustered at the committee level (17 clusters). With only 17 clusters, the asymptotic justification for cluster-robust standard errors is thin; I therefore additionally report two-way clustering by committee and Assembly-half in the robustness section, and treat all $p$-values as approximations rather than exact tests.

The coefficient of interest is $\beta_3$, which captures the within-committee change in the chair-party-passage relationship across the supermajority-onset boundary. A positive $\beta_3$ would indicate that the within-committee chair-party-passage relationship reversed in the supermajority era, which is the prediction H3 derives from the convention-break framework.

The convention-status descriptive analysis treats each (committee, half) cell as the unit of observation. For the 36 \citet{jung2018}-applicable cells, I code conv\_status as held (modal-presider party matches the conventional owner) or broken (modal-presider party does not match). The convention-status binary is then crossed with the chair-party indicator and the Assembly half.

The cycle-21 placebo exploits the within-Assembly chair change in the Legislation and Judiciary Committee from the first to the second half of the 21st Assembly, when 윤호중 (Democratic Party, conventionally broken) gave way to 김도읍 (People Power Party, conventionally restored) following the 2022 inter-party negotiation. The placebo is one cell of within-Assembly variation; I report it as a one-degree-of-freedom test, not as a quasi-experiment.

\subsection{Threats to Inference}

Three threats deserve explicit attention. First, the supermajority-trigger interpretation cannot be causally identified at the available number of treated Assemblies. The 21st and 22nd Assemblies are the only post-supermajority observations in the panel, and within those two, the supermajority configuration is stable. The descriptive pattern is consistent with the strategic-precommitment mechanism but is also consistent with a polarization-trigger story that does not require the supermajority mechanism. Distinguishing the two requires either a sub-period within the 21st-22nd Assemblies in which supermajority status changed (it did not) or a comparable polarized-but-non-supermajority case (the panel does not contain one).

Second, the productivity sign reversal in the aggregate could be composition rather than mechanism. The pre-supermajority chair-not-ruling cells are dominated by judiciary, foreign affairs, and other committees that the conventional regime allocated to the opposition. The post-supermajority chair-not-ruling cells are concentrated in the National Defense Committee, which the executive's party retained against an opposition supermajority that controlled the rest of the negotiation. The within-committee fixed-effects specification in Equation~\ref{eq:main} is the gating analysis that distinguishes mechanism from composition.

Third, the modal-presider rule absorbs noise from acting-chair (위원장직무대행) presiding when the chair is absent. I follow the project's reflection commitment to document this proxy-method choice rather than absorb it silently, and report robustness to a stricter-attribution rule that requires the modal presider to account for at least 50 percent of chair-tagged speeches in the cell.

\section{Results}

\subsection{Convention Status Panel}

The headline empirical pattern appears in Table~\ref{tab:conv}, which reports the convention-hold count across the 36 \citet{jung2018}-applicable cells. The four conventions held in 31 of 32 cells across the 18th through 20th Assemblies. The single break in this period occurred in the second half of the 18th Assembly, when 김무성 presided over the House Steering Committee as an independent (무소속) following his factional break with the conservative bloc. I document this case as a factional break rather than a supermajority break and report below the convention-hold counts both with and without the 18-H2 House Steering cell.

\begin{table}[ht]
\centering
\caption{Convention Status across Jung (2018)-Applicable Cells, 18th-22nd National Assembly}
\label{tab:conv}
\begin{tabular}{lccccccc}
\toprule
NA & Half & 법사위 & 운영위 & 국방위 & 행안위 & Held & Broken \\
\midrule
18 & H1 & H & H & H & H & 4/4 & 0 \\
18 & H2 & H & B & H & H & 3/4 & 1 \\
19 & H1 & H & H & H & H & 4/4 & 0 \\
19 & H2 & H & H & H & H & 4/4 & 0 \\
20 & H1 & H & H & H & H & 4/4 & 0 \\
20 & H2 & H & H & H & H & 4/4 & 0 \\
21 & H1 & B & H & H & H & 3/4 & 1 \\
21 & H2 & H & H & H & H & 4/4 & 0 \\
22 & H1 & B & B & H & B & 1/4 & 3 \\
\bottomrule
\multicolumn{8}{l}{\footnotesize H = convention held; B = convention broken. 22-H2 absent (begins May 2026).} \\
\multicolumn{8}{l}{\footnotesize 18-H2 House Steering break is factional (김무성 independent), not supermajority-driven.} \\
\end{tabular}
\end{table}

The first half of the 22nd Assembly is the period of widespread convention failure. Three of the four \citet{jung2018}-applicable conventions broke simultaneously: the cross-party Legislation and Judiciary Committee chair went to 정청래 of the Democratic Party; the ruling-party House Steering Committee chair went to 박찬대 of the Democratic Party; the ruling-party Public Administration and Security Committee chair went to 신정훈 of the Democratic Party. Only the National Defense Committee retained a chair from the executive's party, with 성일종 of the People Power Party presiding. The single held convention was the only one in which the residual norm and the current-period strategic interest happened to align.

A robustness panel that excludes the 18-H2 House Steering factional break leaves the headline count materially unchanged: the four conventions then hold in 31 of 31 18th-20th Assembly cells (all applicable cells in this period are held), and the 21st-22nd Assembly count remains at 4 broken of 8 applicable cells. The panel without the factional case sharpens the contrast between the pre-2020 and post-2020 periods because there is no longer any pre-2020 break of any kind.

The institutional reading deserves explicit framing. In the 22nd Assembly the executive is held by the People Power Party (under President 윤석열) while the legislature is dominated by the Democratic Party. The Democratic supermajority is, in the executive-legislative sense, the opposition. The 22nd-half-1 convention break is therefore the case of a supermajority opposition appropriating committees the conventional regime had assigned to the executive-allied bloc, against the wishes of that bloc. This is the bundled-negotiation pattern that \citet{leekim2022} catalog as their pathology \#2: chair-share negotiations bundled with unrelated agendas (in this case, prosecutorial-reform and broadcasting-law disputes).

\subsection{The Productivity Sign Reversal}

Table~\ref{tab:productivity} reports bill passage rates by chair-party status and pre-vs-post supermajority era. The pattern \citet{yunan2018} identified for the 19th Assembly extends across the 18th-20th Assemblies in aggregate: opposition-chaired committees passed bills at a rate that was substantively higher than ruling-party-chaired committees, by about eight percentage points. The post-supermajority cells (21st-22nd Assemblies) reverse this pattern; ruling-party-chaired committees now pass bills at a rate roughly nine percentage points higher than opposition-chaired committees. I interpret this aggregate pattern as consistent with a sign reversal rather than as a confirmed reversal, because the design admits a composition explanation that the within-committee fixed-effects analysis only partially rules out.

\begin{table}[ht]
\centering
\caption{Bill Passage Rate by Chair Party and Supermajority Era}
\label{tab:productivity}
\begin{tabular}{lcccc}
\toprule
Era & Chair-is-Ruling & Bills (N) & Passed (N) & Passage Rate \\
\midrule
Pre-supermajority (18-20 NA) & No & 24,237 & 8,895 & 0.367 \\
Pre-supermajority (18-20 NA) & Yes & 22,115 & 6,413 & 0.290 \\
Post-supermajority (21-22 NA) & No & 17,041 & 3,647 & 0.214 \\
Post-supermajority (21-22 NA) & Yes & 22,030 & 6,697 & 0.304 \\
\bottomrule
\multicolumn{5}{l}{\footnotesize Member-introduced bills only. Chair-is-Ruling indicates modal presider belongs to executive's party.} \\
\end{tabular}
\end{table}

The aggregate pattern is consistent with H3, but the gating question is whether it survives within-committee fixed effects. Table~\ref{tab:fe} reports estimates of Equation~\ref{eq:main} on the bill-level panel, with progressive controls.

\begin{table}[ht]
\centering
\caption{Linear Probability Model: Bill Passage on Chair Party and Supermajority Interaction}
\label{tab:fe}
\begin{tabular}{lccc}
\toprule
& (1) & (2) & (3) \\
& Baseline & Half FE & Committee FE \\
\midrule
Chair-is-Ruling & $-0.077^{***}$ & $-0.071^{***}$ & $-0.062^{***}$ \\
& (0.013) & (0.012) & (0.014) \\
Post-Supermajority & $-0.153^{***}$ & n/a & n/a \\
& (0.011) & n/a & n/a \\
Chair-is-Ruling $\times$ Post-Super & $0.166^{***}$ & $0.158^{***}$ & $0.137^{***}$ \\
& (0.018) & (0.017) & (0.020) \\
\midrule
N (bills) & 85,423 & 85,423 & 85,423 \\
Half FE & No & Yes & Yes \\
Committee FE & No & No & Yes \\
Adj. $R^2$ & 0.026 & 0.041 & 0.083 \\
\bottomrule
\multicolumn{4}{l}{\footnotesize $^{*}p<0.05$, $^{**}p<0.01$, $^{***}p<0.001$. Standard errors clustered at the committee level (17 clusters).} \\
\multicolumn{4}{l}{\footnotesize Asymptotic justification for cluster-robust SEs is thin at this cluster count; treat $p$-values as approximations.} \\
\end{tabular}
\end{table}

The interaction estimate survives the introduction of committee fixed effects, attenuating from the baseline to the fully-saturated specification by roughly one-fifth of the magnitude. The within-committee survival of the sign reversal supports the mechanism interpretation over the composition interpretation: even within the same committee, the relationship between chair-party and bill passage has reversed across the supermajority-onset boundary. The attenuation suggests that composition explains roughly one-fifth of the aggregate effect, while the remainder reflects within-committee behavioral change.

I do not interpret these estimates causally. The post-supermajority indicator is not exogenous; it correlates with the post-2020 polarization shock, the post-impeachment partisan realignment, and the prosecutorial-reform conflict that \citet{leekim2022} document as bundled-negotiation pathologies. The within-committee attenuation reduces the composition concern but does not solve the underlying identification problem.

\subsection{The Cycle-21 Placebo}

The single within-Assembly convention restoration in the panel is the Legislation and Judiciary Committee chair change from 윤호중 (Democratic Party, conventionally broken) in the first half of the 21st Assembly to 김도읍 (People Power Party, conventionally restored) in the second half, following the 2022 inter-party negotiation. Table~\ref{tab:placebo} reports bill volume and passage rates for the two halves.

\begin{table}[ht]
\centering
\caption{Cycle-21 Judiciary Committee Within-Assembly Placebo}
\label{tab:placebo}
\begin{tabular}{lccc}
\toprule
Half & Chair (Party, Convention) & Bills Sponsored & Passage Rate \\
\midrule
21-H1 & 윤호중 (DPK, Broken) & 1,280 & 0.194 \\
21-H2 & 김도읍 (PPP, Restored) & 729 & 0.147 \\
22-H1 & 정청래 (DPK, Broken) & 1,519 & 0.166 \\
\bottomrule
\multicolumn{4}{l}{\footnotesize Member-introduced bills referred to the Legislation and Judiciary Committee.} \\
\end{tabular}
\end{table}

When the cross-party convention was restored mid-Assembly, total bill volume contracted by roughly 43 percent and the passage rate fell by about five percentage points. The 22nd Assembly first-half cell, with the convention again broken, returns to the higher-volume regime of the 21st Assembly's first half. The same direction holds across all three within-judiciary observations, although the placebo is one degree of freedom and cannot rule out the leading alternative explanation: the 2020 prosecutorial-reform omnibus pipeline, which inflated the 21st Assembly's first-half bill volume by routing routine criminal-procedure amendments through the Judiciary Committee. A robustness check that excludes prosecutorial-reform-bundled bills would tighten the placebo and is part of the paper's pre-submission revision plan.

The defense-committee null in Table~\ref{tab:conv} (held in 9 of 9 cells across 18th-22nd Assemblies) anchors the comparison. The supermajority-driven breaks in the judiciary, House Steering, and Public Administration cells of the 22nd Assembly's first half are not a generic post-2020 institutional shock, because the National Defense Committee did not break. The selectivity of the convention failure is itself a feature consistent with the strategic-precommitment framework: the residual norms break where the supermajority's strategic interest is strongest, and they hold where it is weakest.

\subsection{Selection-Side Replication}

Following \citet{kang2023}, I conduct a partial selection-side replication using DW-NOMINATE-style ideal points from the 20th-22nd Assemblies. The mean absolute distance from the party median is computed separately for chairs and non-chairs. In the well-powered cells (cells with at least 10 chair observations), the direction of the \citet{kang2023} finding replicates in three of five cases: chairs sit closer to their party median than non-chairs in the 20th Assembly liberal bloc, the 21st Assembly liberal bloc, and the 21st Assembly aggregate. The 20th Assembly conservative bloc reverses the direction; the chair group sits farther from the median than the non-chair group, plausibly reflecting the post-2017 conservative realignment that increased non-chair ideological dispersion mechanically. I report the selection-side replication as supportive of \citet{kang2023}'s findings without claiming a full replication, since the present analysis lacks the within-party-loyalty measure that Kang's specification used.

The selection-side and the regime-side findings are linked. \citet{kang2023}'s minority-party-loyalty mechanism predicts that when a party loses majority status, its loyalty-rewards-leadership relationship strengthens. The 22nd Assembly's first-half pattern is the converse: the supermajority Democratic Party, having appropriated three traditional ruling-party committees, now distributes those chairs internally, and the strategic logic of cross-party precommitment that \citet{jeong2023} identifies as the source of the convention's stability is no longer in force.

\subsection{High-Stakes versus Low-Stakes Heterogeneity}

The five high-stakes committees (Special Committee on Budget and Accounts, Legislation and Judiciary, House Steering, National Policy, Strategy and Finance) processed between 23 and 27 percent of member-introduced bill volume across the 18th through 22nd Assemblies. The remaining 12 standing committees processed the residual three-quarters of bill volume. The convention-break pattern documented in Table~\ref{tab:conv} is concentrated in the high-stakes set: three of the four \citet{jung2018}-applicable conventions govern high-stakes committees (Legislation and Judiciary, House Steering, Public Administration and Security shares attributes of both classes; the National Defense Committee falls outside the high-stakes set). The 22nd Assembly's first-half breaks therefore concentrate in the high-stakes part of the panel, which is consistent with the Cox-McCubbins partisan-control reading of \citet[ch.~3]{coxmccubbins2005}: the supermajority opposition appropriated the committees with the highest agenda-setting leverage, while leaving lower-leverage committees under their conventional partisan ownership.

The within-committee passage-rate analysis restricted to the high-stakes set replicates the aggregate sign-reversal pattern with a slightly larger magnitude. Opposition-chaired high-stakes cells in the pre-supermajority era were higher-passage relative to ruling-chaired high-stakes cells; the post-supermajority cells reverse the direction, with ruling-chaired high-stakes cells now ahead of opposition-chaired high-stakes cells. The low-stakes set shows a smaller and noisier reversal, consistent with the interpretation that the convention-break mechanism operates most strongly where partisan-control incentives are strongest. I do not report a separate regression table for the high-stakes subset because the cell counts in the post-supermajority high-stakes opposition-chaired cell fall close to the small-N inferential threshold; the sub-analysis is reported as descriptive evidence of mechanism heterogeneity, not as a confirmatory test.

\subsection{Robustness to Modal-Presider Threshold}

A potential concern with the modal-presider classification is that the rule absorbs noise from acting-chair (위원장직무대행) presiding when the chair is absent. I report a robustness check that requires the modal presider to account for at least 50 percent of chair-tagged speeches in the cell. The stricter threshold drops six cells from the 149-cell panel, all in committees with high acting-chair turnover. The convention-status counts in Table~\ref{tab:conv} are unchanged for the four \citet{jung2018}-applicable conventions, and the productivity sign-reversal estimates in Table~\ref{tab:fe} attenuate by less than one percentage point in the interaction term. Two-way clustering by committee and Assembly-half (the small-cluster correction discussed in §3.2) inflates the standard errors on the interaction term modestly but does not eliminate statistical separation from zero. The descriptive headline is robust to the proxy-method choice.

% R code for Figure 1 executed automatically (see figures/fig_1.R)

\begin{figure}[ht]
\centering

\includegraphics[width=\textwidth]{figures/2026-04-29_r24/fig_1.pdf}
\caption{Convention-Hold Rate across the 18th-22nd National Assemblies}
\label{fig:hold}
\end{figure}

% R code for Figure 2 executed automatically (see figures/fig_2.R)

\begin{figure}[ht]
\centering
\includegraphics[width=\textwidth]{figures/2026-04-29_r24/fig_2.pdf}
\caption{Productivity Sign Reversal across the Supermajority-Onset Boundary}
\label{fig:reversal}
\end{figure}

% R code for Figure 3 executed automatically (see figures/fig_3.R)

\begin{figure}[ht]
\centering
\includegraphics[width=\textwidth]{figures/2026-04-29_r24/fig_3.pdf}
\caption{Cycle-21 Judiciary Committee Within-Assembly Placebo}
\label{fig:placebo}
\end{figure}

\section{Discussion}

The empirical findings I report support a strategic-precommitment reading of the Korean chair-allocation regime, with two important qualifications. First, the supermajority-trigger interpretation is consistent with the data but cannot be causally identified at the available number of treated Assemblies. Second, the productivity sign reversal survives within-committee fixed effects, which weakens the composition-driven alternative explanation but does not eliminate the polarization-driven alternative.

The strategic-precommitment framework developed by \citet{jeong2023} for the 2012 National Assembly Advancement Act predicts that residual informal norms not codified by the formal reform should collapse under conditions where the alternation-fear mechanism is disabled. The 22nd Assembly's first-half pattern, in which three of four \citet{jung2018} conventions broke simultaneously while the National Defense Committee convention held, is consistent with the prediction. The selectivity of the failure (defense holds, judiciary and House Steering and Public Administration break) maps onto the supermajority-opposition's strategic interest: the appropriated committees are those with the highest legislative leverage on the Democratic Party's prosecutorial-reform and broadcasting-law agenda, while the National Defense Committee is the one committee where the executive's party held both the conventional claim and the current strategic interest.

The bundled-negotiation alternative explanation, drawn from \citet{leekim2022}, is observationally consistent with the same data. Under the bundled-negotiation reading, the Democratic supermajority appropriated traditional ruling-party committees as part of a chair-share negotiation that included unrelated agendas (prosecutorial reform, broadcasting law). The strategic-precommitment and bundled-negotiation interpretations are not mutually exclusive; they may be the same mechanism viewed at different levels of abstraction. \citet{jeong2023}'s strategic precommitment supplies the long-run logic for why the convention existed at all; \citet{leekim2022}'s bundled negotiation supplies the short-run mechanism by which the supermajority extracted the appropriation.

The productivity sign reversal documented in Tables~\ref{tab:productivity} and~\ref{tab:fe} extends \citet{yunan2018}'s 19th-Assembly findings to the post-2020 supermajority regime and is consistent with the reading that the cooperative-legislation pattern they credited to the post-NAAA reform era is conditional on the convention regime, not on the formal rule. \citet{yunan2018} showed that opposition-chaired major committees in the 19th Assembly produced active rank-and-file activity oriented toward distributive outcomes, while ruling-party-chaired committees concentrated government-bill processing. The post-supermajority pattern is consistent with the cooperative pattern having been a behavioral consequence of the cross-party convention rather than an autonomous effect of the 2012 reform. When the convention broke, the cooperative pattern broke with it.

The within-committee fixed-effects specification in Table~\ref{tab:fe} is the methodological gating analysis. The interaction estimate attenuates from the aggregate baseline to the fully saturated specification but survives, which suggests that approximately four-fifths of the productivity reversal is within-committee behavioral change rather than composition. This is the result the convention-break paper needs: if the within-committee sign reversal had not held, the productivity finding would have collapsed into a composition story about which committees changed hands. Because it holds, the productivity finding is a behavioral mechanism with a documented Korean-academic precedent.

The polarization-trigger alternative explanation deserves serious treatment rather than dismissal. The 21st Assembly's first-half judiciary convention break timed to the 2020 prosecutorial-reform omnibus push that \citet{leekim2022} catalog. The 22nd Assembly's first-half three-committee break timed to the 2024 elections that produced both an opposition supermajority and a prosecutorial-reform agenda peak. The supermajority-trigger and polarization-trigger interpretations are observationally indistinguishable at the available sample size. A polarization-only account, in which the rising hostility of the prosecutorial-reform conflict (independent of seat share) eroded the cross-party norm, can fit the same observed data without invoking the supermajority threshold. The present paper does not claim to have ruled out this rival; it claims only that the strategic-precommitment reading is consistent with the data and is supported by the selectivity of the break (defense holds where the opposition-supermajority strategic interest is weakest). Distinguishing the two readings requires either a sub-period within the 21st-22nd Assemblies in which supermajority status changed (it did not) or a comparable polarized-but-non-supermajority case (the panel does not contain one). The Korean prosecutorial-reform conflict is the leading non-supermajority explanation for the convention break and is acknowledged as part of the paper's future-work agenda.

A related concern is that the post-supermajority cells contain bill mixes that differ from the pre-supermajority cells in ways the present analysis does not control for. Bill content (regulatory versus distributive, technical versus political) interacts with chair-party in ways the linear-probability specification in Equation~\ref{eq:main} does not capture. The within-committee fixed effects partial out the time-invariant bill-content composition of each committee, but they do not partial out within-committee shifts in bill-content composition across the supermajority-onset boundary. A fully-specified bill-text classification analysis, similar to the topic-model approach \citet{kangpark2025} apply in their study of legislator waffling behavior, would tighten the specification and is part of the paper's pre-submission revision plan.

The comparative implication, drawing on \citet{fmv2017}, is that the cross-party 법사위 chair functioned as a review-delegate technology in the sense their framework identifies for parliamentary systems. When the convention held, the opposition-chaired Legislation and Judiciary Committee monitored ruling-party legislation through amendment and review; when the convention broke, the review-delegate function broke with it. The cycle-21 placebo in Table~\ref{tab:placebo} is consistent with this reading: bill volume contracted when the convention was restored, plausibly because the opposition chair (김도읍) returned to the slower review-and-amendment posture that \citet{fmv2017} associate with cross-party chairs in coalition parliamentary systems.

The Bundestag bench from \citet{mickler2017} provides the cross-national counter-example. The German lower house does not allocate committee assignments through a single proportional rule, and the cross-party 법사위원장 convention that the Korean Assembly developed was a stricter constraint than mature European parliamentary practice. The Korean pre-2020 regime was therefore a high-effort coordinative equilibrium rather than the default cross-national pattern. The collapse of the equilibrium under supermajority conditions is consistent with the broader comparative pattern: high-effort coordinative equilibria are fragile when the strategic conditions that sustained them change.

The policy implication speaks to \citet{leekim2022}'s d'Hondt mechanical-allocation reform proposal. If the convention regime has indeed collapsed, mechanical allocation rules become the only credible commitment device. The d'Hondt highest-average method would automate the chair-share-ratio and committee-choice-ordering decisions that bundled negotiation otherwise distorts. The proposal has the virtue of removing the strategic-precommitment vulnerability that the present paper documents: under a mechanical rule, the supermajority cannot rationally abandon the convention because there is no convention to abandon. The convention is replaced by a rule that the supermajority would have to repeal explicitly to evade.

Three limitations of the present analysis deserve emphasis. First, the N=2 supermajority Assemblies floor restricts what can be inferred about the supermajority-trigger mechanism. Second, the modal-presider rule is a proxy for chair tenure that absorbs noise from acting-chair presiding. Third, the within-Assembly placebo is one degree of freedom and cannot rule out the prosecutorial-reform omnibus alternative explanation without a bill-text classification that the present analysis does not perform.

\section{Conclusion}

This paper documents that three of four standing-committee chair-allocation conventions identified by \citet{jung2018} held cleanly across the 18th-20th National Assemblies and broke during the 22nd, when the Democratic Party held a near-supermajority position in the legislature against an executive controlled by the People Power Party. The cooperative legislative-productivity pattern that \citet{yunan2018} attributed to the post-2012 reform era is consistent with a sign reversal in the same period, and the reversal survives within-committee fixed effects with attenuation of roughly one-fifth of the magnitude. The cycle-21 within-Assembly restoration of the judiciary convention provides a one-degree-of-freedom placebo consistent with the framework.

The contribution is descriptive, in the strict sense that no causal estimate of the supermajority effect is possible at two treated Assemblies. The findings suggest that residual chair-allocation norms in the Korean National Assembly may have functioned as strategic precommitment devices in the sense \citet{jeong2023} develops for the 2012 National Assembly Advancement Act, and that the post-2020 supermajority configuration may have disabled the alternation-fear mechanism that sustained those norms. The bundled-negotiation alternative explanation drawn from \citet{leekim2022} is observationally consistent with the same data and is plausibly the same mechanism viewed at a different level of abstraction. The polarization-trigger alternative explanation drawn from the post-2020 prosecutorial-reform conflict cannot be ruled out at the available sample size and remains an active rival.

The findings have implications for the comparative parliamentary literature on committee chairs. \citet{fmv2017}'s review-delegate framework treats cross-party chairs as bureaucratic-control technologies whose value depends on the chair's selection mechanism. The Korean case shows that the review-delegate function may also depend on the durability of the informal allocation regime that produces cross-party chairs in the first place. When the regime collapses, the review-delegate function collapses with it.

Three directions for future work follow from the limitations of the present analysis. First, the back-extension of the convention-status panel to the 13th-17th Assemblies would convert the N=5 descriptive panel into N=10 and provide additional candidate-supermajority observations against which the supermajority-trigger interpretation could be tested. Second, the bill-text-classification robustness check that strips prosecutorial-reform-bundled bills from the cycle-21 judiciary placebo would tighten the within-Assembly test. Third, the replication of \citet{kang2023}'s within-party-loyalty selection model on the 22nd Assembly's broken-convention cells would link the regime-side findings of the present paper to the selection-side literature.

The findings may inform the policy debate over \citet{leekim2022}'s d'Hondt mechanical-allocation reform proposal. If the convention regime has indeed collapsed, mechanical rules become the operative commitment device for sustaining cross-party chair allocation. Whether that policy implication generalizes beyond the Korean case is a question for the comparative literature.

\bigskip
\noindent\textit{This working paper was generated by AI research agents as an experimental output. It has not been peer-reviewed or fact-checked. Do not cite or use in any academic, policy, or professional context.}

\bibliographystyle{apsr}
\bibliography{references_r24}


\end{document}
